\section*{Chapter \ref{ch:g6:school-and-work} --- School and Work}

\begin{solution}{exo:g6:school-and-work:1}
About two in three ($65\,\%$) in 1975 and almost nine in ten ($88\,\%$) in 2024;
the table gives $65.4\,\%$ and $87.7\,\%$.
\end{solution}

\begin{solution}{exo:g6:school-and-work:2}
South Korea, with about 1\,900 hours (the table gives $1\,910$).
\end{solution}

\begin{solution}{exo:g6:school-and-work:3}
$1\,800 \div 45 = 40$ hours a week.
\end{solution}

\begin{solution}{exo:g6:school-and-work:4}
About half ($50\,\%$): the farming layer reaches about the 50 gridline in 2000.
\end{solution}

\begin{solution}{exo:g6:school-and-work:5}
(a) The change is $2\,427 - 1\,335 = 1\,092$ hours, and $1\,092 \div 2\,427
\approx 0.450$: a fall of $45.0\,\%$. (b) $1\,092 \div 52 = 21$ hours fewer per
week.
\end{solution}

\begin{solution}{exo:g6:school-and-work:6}
(a) Farming $12 \times 0.35 = 4.2$ million; industry $12 \times 0.25 = 3$ million;
services $12 - 4.2 - 3 = 4.8$ million. (b) $100 - 35 - 25 = 40\,\%$.
\end{solution}

\begin{solution}{exo:g6:school-and-work:7}
In the early 2010s (the table gives 2011): from then on, the services layer at the
top is thicker than the farming layer at the bottom.
\end{solution}

\begin{solution}{exo:g6:school-and-work:8}
(a) South Korea's line rises from 1953 to 1969, while the country was growing
richer fast; Japan's rises from 1950 to 1960 too. (b) ``In each of the six
countries, workers worked fewer hours in 2023 than in 1960, when they were much
poorer.''
\end{solution}

\begin{solution}{exo:g6:school-and-work:9}
(a) $4.5 \div (4.5 + 5.5) = 4.5 \div 10 = 45\,\%$. (b) No. It says what share of
\emph{workers} are women. To know what share of \emph{women} work, one needs the
number of women of working age, for example 9 million: then $4.5 \div 9 = 50\,\%$.
\end{solution}

\begin{solution}{exo:g6:school-and-work:10}
(a) $48.9 - 51.2$: a fall of $2.3$ percentage points. (b) $2.3 \div 51.2 \approx
0.045$: a fall of $4.5\,\%$. (c) For example: young women stay at school longer
(check: the share of women aged 15 to 24 in school, at both dates); in some large
countries fewer women work on family farms as farming shrinks (check: the share,
country by country, and the share of farm jobs). Any explanation must be checked
country by country, because the world figure mixes rises and falls.
\end{solution}

\begin{solution}{exo:g6:school-and-work:11}
(a) 1990: $600\,000 \times 0.60 = 360\,000$; 2020: $900\,000 \times 0.25 =
225\,000$. (b) 1990: $600\,000 - 360\,000 = 240\,000$; 2020: $900\,000 - 225\,000 =
675\,000$. (c) Faster: all jobs rose by $300\,000 \div 600\,000 = 50\,\%$, jobs
outside farming by $435\,000 \div 240\,000 \approx 181\,\%$.
\end{solution}

\begin{solution}{exo:g6:school-and-work:12}
(a) For example: safer work (fewer accidents and deaths at work); less heavy
physical effort; more pleasant places to work; work closer to home. (b) For
example: more stress, more evening or night work, less secure contracts.
(c) For example, for safety: the number of deaths and injuries at work per
100\,000 workers, year after year.
\end{solution}

\begin{solution}{pb:g6:school-and-work:1}
\textbf{1.} Grandfather $10 \times 6 \times 50 = 3\,000$ hours; mother $8.5 \times 5
\times 47 = 1\,997.5$ hours; grandson $7 \times 5 \times 46 = 1\,610$ hours.
\textbf{2.} $(3\,000 - 1\,610) \div 3\,000 = 1\,390 \div 3\,000 \approx 0.463$: a fall
of $46.3\,\%$.
\textbf{3.} The grandfather worked $10 \times 6 = 60$ hours a week: $60 - 35 = 25$
hours more.
\textbf{4.} 1965: $400\,000 \times 0.50 = 200\,000$; 2025: $600\,000 \times 0.08 =
48\,000$.
\textbf{5.} 1965: $400\,000 \times 0.20 = 80\,000$; 2025: $600\,000 \times 0.70 =
420\,000$.
\textbf{6.} No. Each farm worker may produce much more food, with machines and
better seeds; and Brenland may buy food from abroad. One needs the quantity of
food produced, and what Brenland buys and sells abroad.
\textbf{7.} Grandfather $60 \div 0.50 = 120$ hours; mother $180 \div 6 = 30$ hours;
grandson $300 \div 20 = 15$ hours.
\textbf{8.} Grandfather $3\,000 \div 120 = 25.0$; mother $1\,997.5 \div 30 \approx
66.6$; grandson $1\,610 \div 15 \approx 107.3$ bicycles.
\textbf{9.} No: it rose from 25 to $107.3$. The wage per hour rose much faster than
the price of a bicycle, so far more than the lost hours were made up.
\textbf{10.} No. The bicycle's price was 3 times as high ($180 \div 60$), so an hour
of her work bought $12 \div 3 = 4$ times as much bicycle (30 hours instead of 120),
and a year of her work, with fewer hours, bought $66.6 \div 25 \approx 2.7$ times
as many bicycles. A wage must be compared with prices.
\textbf{11.} For example: how hard, dangerous or tiring the work was; whether the
job was secure; the taxes paid; the prices of food and housing, which may have
changed differently from the price of a bicycle.
\end{solution}
